% !TEX root = ../projekt.tex
% =============================================================================
% STAV: rozpracované, dopísať po vyhodnotení variantov (bod 6 zadania).
% =============================================================================
\chapter{Záver}
\label{chap:zaver}

Práca sa zaoberala distribúciou paketov do DMA kanálov jednotkou RSS frameworku NDK. Vznikol
simulátor jednotky RSS pre najviac 128 kanálov, ktorý pracuje so záznamami reálnej prevádzky
a záťaž hodnotí preťažením, Jainovým indexom a normalizovaným $\chi^2$.

Analýza vzorky \texttt{odid1} ukázala, že nerovnomernú záťaž nespôsobuje kvalita hashu, ale
rozdielna veľkosť tokov. Medzi piatimi hashovacími funkciami sa nepreukázal rozdiel väčší než
vplyv náhodného kľúča, slabšia funkcia sa však môže prejaviť pri vstupe s~nízkou entropiou. Referenčné orákulum, ktoré zachováva afinitu toku, znížilo preťaženie
oproti Toeplitzu o~63 až 88\,\% pri váhach packet a byte, čo ukazuje veľkú rezervu v~spôsobe
výberu kanála.

\TODO{Doplniť výsledky navrhnutých variantov a odporúčanie pre NDK.}

\paragraph{Možnosti pokračovania.}
\TODO{Napríklad implementácia najlepšej varianty vo VHDL, meranie na reálnom hardvéri,
vyhodnotenie na paketových stopách namiesto tokových záznamov.}
